\chapter{Setting Up Your Development Environment}
\label{ch:firstbuild}

This chapter helps you get a built, runnable \prg on Windows,
Linux, or macOS. The instructions follow the platform's own setup guides
\citep{prg32repo}. The common foundation on every operating system is the same:
install Espressif's ESP-IDF, the development framework for these \riscv
microcontrollers, then build the firmware for your chosen target. Pick your
operating system's section, then read the shared sections on QEMU and the build
commands.

\section{The one idea behind every setup}

Whatever your operating system, the toolchain is ESP-IDF version 5.3 or newer,
configured for the \code{esp32c3} and \code{esp32c6} targets: the first for the
QEMU emulator, the second for the physical board \citep{prg32repo,prg32qemu}. Once
ESP-IDF is installed and "sourced" (made active in your shell), the command
\code{idf.py} builds, flashes, and runs everything. Almost every setup problem
reduces to ESP-IDF not being active in the current shell.

\begin{prgwarn}
ESP-IDF must be activated in \emph{each new terminal} before \code{idf.py} works.
If you see \code{idf.py: command not found} or a missing \code{riscv32-esp-elf-gcc},
ESP-IDF is simply not sourced in that shell. This is the single most common
setup error, and the fix is always the same: re-run the export step
\citep{prg32repo}.
\end{prgwarn}

\section{Environment installation}

Install ESP-IDF and any platform prerequisites first. After ESP-IDF is active,
\code{idf.py} becomes available in the shell and the rest of this workflow can
run.

\subsection{Windows}

The platform recommends a classroom path built on Visual Studio Code with the
Espressif ESP-IDF extension, which installs the toolchain for you, plus the
PlatformIO and Microsoft C/C++ and Python extensions \citep{prg32repo}. There is
also a standalone path using the official ESP-IDF Tools Installer, after which you
work from the "ESP-IDF PowerShell" shortcut it creates.

\lstinputlisting[
    style=bash,
    caption={%\fname{windows\_setup.sh} \protect\
       Windows standalone path, run from the ESP-IDF PowerShell shortcut.}
]{scripts/setup/windows_setup.sh}

\begin{prgwarn}
On Windows, use the ESP-IDF PowerShell or ESP-IDF Command Prompt, not a plain
terminal where \code{idf.py} has not been exported. If flashing fails, check Device
Manager for the board's serial port and pass it explicitly with \code{-p COMx}; if
the monitor cannot open the port, close every other program using it (Arduino
Serial Monitor, other ESP-IDF monitors, and so on) \citep{prg32repo}.
\end{prgwarn}

\subsection{Linux}

On Debian or Ubuntu, install the build prerequisites with the system package
manager, then ESP-IDF as above \citep{prg32repo}. Other distributions use their own
package names for the equivalent libraries.

\lstinputlisting[
    style=bash,
    caption={%\fname{linux\_prereqs.sh} \protect\
       Linux (Debian/Ubuntu) prerequisites, then ESP-IDF.}
]{scripts/setup/linux_prereqs.sh}

\subsection{macOS}

The platform's macOS path, tested on Apple Silicon, installs the prerequisites with
Homebrew and then ESP-IDF \citep{prg32repo}.

\lstinputlisting[
    style=bash,
    caption={%\fname{macos\_setup.sh} \protect\
       macOS setup with Homebrew and ESP-IDF.}
]{scripts/setup/macos_setup.sh}

A convenient habit is an alias so a single word activates ESP-IDF in any shell
\citep{prg32repo}.

\lstinputlisting[
    style=bash,
    caption={%\fname{alias\_activate\_idf.sh} \protect\
       A quality-of-life alias for activating ESP-IDF.}
]{scripts/setup/alias_activate_idf.sh}

\section{Physical board install}

The physical ESP32\nobreakdash-C6 board uses the ESP-IDF toolchain for the
\code{esp32c6} target. Make sure your board is connected, ESP-IDF is active, and
you have serial-port access configured for your OS.

\subsection{Windows}

Use the ESP-IDF PowerShell or ESP-IDF Command Prompt created by the installer.
Once ESP-IDF is active, the same firmware build and flash commands work.

\subsection{Linux}

To flash a real board, your user must be allowed to use the serial port. Add
yourself to the \code{dialout} group, then log out and back in
\citep{prg32repo}.

\lstinputlisting[
    style=bash,
    caption={%\fname{linux\_serial\_access.sh} \protect\
       Granting serial-port access on Linux.}
]{scripts/setup/linux_serial_access.sh}

\begin{prgaside}
ESP32\nobreakdash-C6 boards usually appear as \fname{/dev/ttyACM0}, while USB-bridge boards may
appear as \fname{/dev/ttyUSB0} \citep{prg32repo}. If flashing fails although the
board is visible, the cause is almost always the \code{dialout} membership not yet
taking effect. Log out and back in after adding yourself to the group.
\end{prgaside}

\subsection{macOS}

On macOS, once ESP-IDF is active the physical board flash workflow is the same as
on the other platforms. Use the native serial device exposed by the board and
flash it from an active ESP-IDF shell.

\subsection{Build and flash}

With ESP-IDF active, build and flash the firmware for the physical board using the
following commands. On the \textbf{physical board}, the resident firmware starts its own Wi-Fi access
point named \code{PRG32}. You connect your computer to that network and upload the
cartridge over HTTP; the firmware stores it and runs it, no reflashing required.

\lstinputlisting[
    style=bash,
    caption={%\fname{flash\_board.sh} \protect\
       Building and flashing for the physical ESP32\nobreakdash-C6 board.}
]{scripts/setup/flash_board.sh}

\section{QEMU install}

Installing QEMU lets you run the PRG32 platform on the desktop using the
\code{esp32c3} target.

Linux and MacOS hosts need a few native libraries for the QEMU display window
\citep{prg32qemu}.

\subsection{macOS}

\lstinputlisting[
    style=bash,
    caption={%
       QEMU window libraries for MacOS.}
]{scripts/setup/qemu_window_libs_macos.sh}

\subsection{Linux}

\lstinputlisting[
    style=bash,
    caption={%
       QEMU window libraries for Linux.}
]{scripts/setup/qemu_window_libs_linux.sh}

\subsection{Install the Risc-V QEMU tool}

Install Espressif's \riscv QEMU after ESP-IDF, then re-activate ESP-IDF so the
new tool is on the path \citep{prg32qemu}.

\lstinputlisting[
    style=bash,
    caption={%\fname{install\_qemu.sh} \protect\
       Installing the RISC-V QEMU tool and re-activating ESP-IDF.}
]{scripts/setup/install_qemu.sh}

\section{Build and run}

You will build every game for two targets. The commands differ only in a build
directory and a configuration profile; the source does not change at all
\citep{prg32qemu}. Make sure that ESP-IDF 
is active by sourcing \code{export.sh} in the \code{esp-idf} directory.

For the \textbf{physical ESP32\nobreakdash-C6 board}, the build uses a directory named
\fname{build-esp32c6} and the board configuration profile. 

\lstinputlisting[
    style=bash,
    caption={%
        \fname{building\_and\_flashing\_esp32c6.sh} \protect\\
        Building and flashing for the physical board (abbreviated; see Chapter~\ref{ch:firstbuild}).%
    }
]{scripts/setup/building_and_flashing_esp32c6.sh}

For the \textbf{QEMU emulator}, the build uses \fname{build-qemu} and the QEMU
profile, and targets the ESP32\nobreakdash-C3 emulator. On \textbf{QEMU}, there is no Wi-Fi to upload over, so you build the cartridge
against the QEMU firmware ELF and stage it directly into the emulator's flash
image. If the flash image does not yet exist, run QEMU once so ESP-IDF creates it,
quit, then stage \citep{prg32qemu}.

\subsection{Windows}

\lstinputlisting[
    style=bash,
    caption={%
        \fname{building\_and\_flashing\_qemu\_windows.ps1} \protect\\
        Building and running for the QEMU emulator on Windows.%
    }
]{scripts/setup/building_and_flashing_qemu_windows.ps1}

\subsection{Linux and MacOS}

\lstinputlisting[
    style=bash,
    caption={%
        \fname{building\_and\_flashing\_qemu\_shell.sh} \protect\\
        Building and running for Linux or MacOS.%
    }
]{scripts/setup/building_and_flashing_qemu_shell.sh}

\begin{prgwarn}
Keep the two build directories strictly separate. Flashing a \fname{build-qemu}
image to a real board produces a black screen, because the emulator build uses a
virtual display panel that does not exist on real hardware. If your board shows
nothing, the first thing to check is that you built in \fname{build-esp32c6}, not
\fname{build-qemu} \citep{prg32qemu}.
\end{prgwarn}

\begin{prgaside}
Why does the emulator target an ESP32\nobreakdash-C3 while the board is an ESP32\nobreakdash-C6? Both
are \riscv microcontrollers from the same family, and both run the same 32-bit
\riscv instructions. Espressif maintains a mature QEMU graphics path for the
ESP32\nobreakdash-C3, so \prg uses it for desktop emulation, while the physical kits are
built around the newer ESP32\nobreakdash-C6 \citep{prg32qemu}. From your game's point of
view the difference is invisible: the calling convention and the framework
interface are identical.
\end{prgaside}

\section{Confirming it works}

A healthy board boot logs the firmware entering \code{app\_main}, the configured
display pins, and finally "PRG32 runtime initialized," then shows the splash
\citep{prg32tutorial}. The platform also ships a smoke test that exercises the whole
workflow: environment, firmware build, cartridge build, and QEMU staging, and
prints \code{=== SMOKE TEST PASSED ===} when all is well \citep{prg32repo}.

\lstinputlisting[
    style=bash,
    caption={%\fname{smoke\_test.sh} \protect\
       Running the end-to-end smoke test.}
]{scripts/setup/smoke_test.sh}

\begin{prgcheck}
Your environment is ready when, on your operating system, \code{idf.py} runs in an
ESP-IDF shell, the QEMU build opens a virtual screen reaching "PRG32 runtime
initialized," and, if you have a board, the firmware flashes and shows its
splash. With that in place, every chapter and every example in this book is
buildable on your machine.
\end{prgcheck}
